%% toc: P15 174 174.5,174.6,174.7,174.8
\chapter{Treatment Versions, Interference, Adherence, Missing Outcomes, and Sample Size}
\label{sec:p15-w02b-p15-076}

A randomized assignment licenses a comparison of assignment cells only
through an observation model: a statement of what is recorded for each unit
once the cells have been assigned. The usual model says that the recorded
response of unit \(i\) is its potential response at its own assigned cell.
That one sentence combines four premises: each cell is delivered in one
version, or in versions that do not matter; a unit's response does not
depend on the cells of other units; the analysis concerns assignment, not
the treatment actually received; and every response is recorded. This
chapter replaces each premise by an explicit alternative, proves what the
assignment comparison then estimates, and proves what further premise
recovers the original target. Its last part turns the precision formulas of
\cref{sec:p15-factorial-uncertainty} into sample-size arguments: how many
units, replicates or blocks reach a stated interval half-width or power for
a specified contrast, under a stated mean and variance model, and what
such an argument does not establish.

The contrast derivations for main effects and interactions, with explicit
coding and the conditional, averaged and difference-of-differences
comparisons, are in
\cref{sec:p15-factorial-two-by-two,sec:p15-factorial-full}. Fractional
designs, defining relations, aliasing and estimability are in
\cref{sec:p15-factorial-estimability,sec:p15-factorial-fractions,sec:p15-factorial-half}.
They are used here without restatement.

\paragraph{Inputs and conclusions.}
The inputs are a finite set of recruited units, a finite set of assignment
cells, a known assignment law with positive marginal cell probabilities,
contrast weights, and one of the observation models defined below. The
conclusions are exact expectations of the weighting statistic, equalities
between estimands under named premises, constructions showing that a
target is not determined by what is recorded, sharp bounds, and minimal
sample sizes under explicit normal linear models.

\paragraph{Sources.}
The dependence of a counterfactual on how a treatment is specified, and
the consequence that one treatment name can denote different interventions
in different studies, are discussed in
\cite[\S3.4, pp.~34--36]{HernanRobins2020WhatIf2026}. Interference, and the
need to name the allocation of the other units when it is present, are
discussed in \cite[Fine Point~1.1, p.~5]{HernanRobins2020WhatIf2026}; cluster
allocation as a protection against contamination, and the unit-of-analysis
error, in \cite[\S\S23.1.1, 23.1.3]{HigginsEldridgeLi2024Cochrane23}. The
compliance types, the monotonicity condition and the identity between the
ratio of assignment effects and the effect in compliers are presented, for
a population, in \cite[\S16.4, pp.~217--218, and Technical Point~16.7,
p.~224]{HernanRobins2020WhatIf2026}; the distinction between the effect of
assignment and the effect of adhering to the protocol is drawn in
\cite[\S22.1, p.~306]{HernanRobins2020WhatIf2026} and
\cite[\S\S8.2.2, 8.4.4]{HigginsSavovicEtAl2024Cochrane08}. Data-only bounds
for a bounded outcome are described in
\cite[Technical Point~16.2, p.~212]{HernanRobins2020WhatIf2026}, and the
dependence of complete-case bias on the relation between missingness and the
true outcome within each arm in
\cite[\S\S8.5.1--8.5.2]{HigginsSavovicEtAl2024Cochrane08}. The design effect
\(1+(M-1)\,\mathrm{ICC}\) for clusters of average size \(M\) is the one
stated and applied in \cite[\S\S23.1.4--23.1.4.1]{HigginsEldridgeLi2024Cochrane23}.
The finite-unit formulations, all theorems, proofs, constructions and worked
numbers below are stated and proved in this handbook; the worked tables are
hypothetical and are not observations from a conducted study.

\section{The standing setting and the weighting statistic}
\label{sec:p15-w02b-p15-076-setting}

\begin{hbconvention}{def:p15-w02b-p15-076-setting}{Assignment law, recorded response and weighting statistic}
Fix \(N\geq1\) recruited units, a nonempty finite set \(\mathcal T\) of
cells, real weights \(w=(w_t)_{t\in\mathcal T}\), and a random assignment
vector \(Z=(Z_1,\ldots,Z_N)\in\mathcal T^N\) with known law
\(\varpi(z)=\Pr(Z=z)\) and marginal probabilities
\(\pi_i(t)=\Pr(Z_i=t)>0\). All random variables are defined on one finite
outcome space. An \emph{observation model} specifies a real recorded
response \(Y_i^{\mathrm{rec}}\) for each unit as a function of \(Z\) and of
any further random quantities the model introduces. The \emph{weighting
statistic} of a recorded response is
\[
 \widehat\tau=\frac1N\sum_{i=1}^N\sum_{t\in\mathcal T}
   \frac{w_t\mathbf 1_{\{Z_i=t\}}Y_i^{\mathrm{rec}}}{\pi_i(t)},
\]
the statistic of \cref{thm:p15-factorial-assignment}. The same formula
applied to another recorded variable, such as a receipt or recording
indicator, is the weighting statistic of that variable.
\index{weighting statistic}\index{observation model}
\end{hbconvention}

Under the observation model of \cref{def:p15-factorial-assignment},
\(Y_i^{\mathrm{rec}}=Y_i(Z_i)\), and the target is
\(\tau=N^{-1}\sum_i\sum_tw_tY_i(t)\). The following lemma holds for every
observation model and is the common first step of each later result.

\begin{hblemma}{thm:p15-w02b-p15-076-weighting}{proved=annex:p15-w02b-p15-076-weighting}{What the weighting statistic estimates under any observation model}
For every observation model in the standing setting,
\[
 \mathbb E[\widehat\tau]=\frac1N\sum_{i=1}^N\sum_{t\in\mathcal T}
   w_t\,\mathbb E\bigl[Y_i^{\mathrm{rec}}\bigm| Z_i=t\bigr].
\]
In particular \(\widehat\tau\) is unbiased for \(\tau\) whenever
\(\mathbb E[Y_i^{\mathrm{rec}}\mid Z_i=t]=Y_i(t)\) for every unit \(i\)
and every cell \(t\) with \(w_t\ne0\).
\end{hblemma}

The lemma uses only that each \(\pi_i(t)\) is positive. It therefore
shifts every question about an observation model to one computation: the
mean recorded response of unit \(i\) over the assignments that place it in
cell \(t\).

\section{Treatment versions and stable versions}
\label{sec:p15-w02b-p15-076-versions}

\begin{hbdefinition}{def:p15-w02b-p15-076-versions}{Versions, a version mechanism and the stable-version premise}
In the standing setting, let each cell \(t\) have a nonempty finite set
\(\mathcal V_t\) of \emph{versions}, and let unit \(i\) have a real
potential response \(Y_i(t,v)\) for every \(t\) and \(v\in\mathcal V_t\).
A \emph{version mechanism} is a random tuple
\(\xi=(\xi_1,\ldots,\xi_N)\) independent of \(Z\), together with maps
\(\phi_i\) such that \(\phi_i(t,\xi_i)\in\mathcal V_t\). The delivered
version is \(V_i=\phi_i(Z_i,\xi_i)\), and \emph{version consistency} is
the observation model
\[
 Y_i^{\mathrm{rec}}=Y_i(Z_i,V_i).
\]
Write \(\lambda_i(v\mid t)=\Pr(\phi_i(t,\xi_i)=v)\) for the version law
of unit \(i\) in cell \(t\). The \emph{version-mixture estimand} is
\[
 \tau_\lambda=\frac1N\sum_i\sum_tw_t\sum_{v\in\mathcal V_t}
   \lambda_i(v\mid t)\,Y_i(t,v).
\]
For a selection \(\mathbf v=(v_t)_{t\in\mathcal T}\) with
\(v_t\in\mathcal V_t\), the \emph{version-specific target} is
\(\theta_{\mathbf v}=N^{-1}\sum_i\sum_tw_tY_i(t,v_t)\).
The \emph{stable-version premise} holds when, for every unit \(i\) and
cell \(t\), \(Y_i(t,\cdot)\) is constant on \(\mathcal V_t\); its common
value is then written \(Y_i(t)\).
\index{treatment versions}\index{stable-version premise}\index{consistency!version}
\end{hbdefinition}

The mechanism may treat units differently, and \(\phi_i\) may depend on
unit characteristics that also affect the response. It reads only the
unit's own cell. A mechanism in which the version delivered to one unit
depends on the cells of other units, for example through a shared
deliverer's workload, is a form of interference and is treated in
\cref{sec:p15-w02b-p15-076-interference}.

\begin{hbtheorem}{thm:p15-w02b-p15-076-versions}{proved=annex:p15-w02b-p15-076-versions}{What the assignment comparison estimates when versions differ}
In the setting of \cref{def:p15-w02b-p15-076-versions}:
\begin{enumerate}
\item \(\mathbb E[\widehat\tau]=\tau_\lambda\) for every version
      mechanism.
\item Under the stable-version premise,
      \(\tau_\lambda=\theta_{\mathbf v}=\tau\) for every version mechanism
      and every selection \(\mathbf v\).
\item If \(Y_i(t,\cdot)\) is not constant on \(\mathcal V_t\) for some unit
      \(i\) and some cell \(t\) with \(w_t\ne0\), then two version mechanisms
      with the same potential responses and the same assignment law give
      different values of \(\tau_\lambda\).
\item Suppose only \((Z,Y_1^{\mathrm{rec}},\ldots,Y_N^{\mathrm{rec}})\) is
      recorded. If, for some unit \(i\) and cell \(t\) with \(w_t\ne0\),
      \(Y_i(t,v_t)\ne Y_i(t,v')\) for some \(v'\in\mathcal V_t\), then a
      second model, obtained by exchanging the two versions \(v_t\) and
      \(v'\) for that unit and cell, has pointwise the same recorded
      vector and a different value of \(\theta_{\mathbf v}\). If
      \(\lambda_i(v_t\mid t)=\lambda_i(v'\mid t)\), the second model also
      has the same version laws.
\item Suppose \(V\) is also recorded. If
      \(\lambda_i(v_t\mid t)>0\) for every unit \(i\) and every cell \(t\)
      with \(w_t\ne0\), then
      \[
       \theta_{\mathbf v}=\frac1N\sum_i\sum_tw_t
        \frac{\mathbb E\bigl[\mathbf 1_{\{Z_i=t,\,V_i=v_t\}}Y_i^{\mathrm{rec}}\bigr]}
             {\Pr(Z_i=t,V_i=v_t)},
      \]
      and the statistic obtained by replacing each expectation with the
      realized indicator times \(Y_i^{\mathrm{rec}}\) is unbiased for
      \(\theta_{\mathbf v}\). If \(\lambda_i(v_t\mid t)=0\) for some such
      unit and cell, changing \(Y_i(t,v_t)\) changes
      \(\theta_{\mathbf v}\) and leaves the joint law of all recorded
      variables unchanged.
\end{enumerate}
\end{hbtheorem}

Part~1 says that an experiment with several versions estimates a contrast
of version mixtures. The mixture weights are set by the delivery process
in force, and that process need not be under the experimenter's control.
Part~3 shows that the estimand then belongs to the pair of potential
responses and delivery process, not to the cell labels alone. Parts~4 and~5
separate two different questions. Recording which version was delivered
does not make the mixture contrast a version-specific one; it permits a
version-specific target only for versions that every unit receives with
positive probability.

\begin{hbexample}{def:p15-w02b-p15-076-versions-example}{A checking prompt delivered in two ways}
Use the aid-present cells \(01\) and \(11\) of
\cref{def:p15-factorial-coding}, with weights \(w_{01}=-1\) and
\(w_{11}=1\), and four participants. Cell \(01\) has one version. The
checking prompt of cell \(11\) is either read aloud, version \(e\), or
printed only, version \(p\). Stipulate the potential numbers of
inconsistent pairs
\[
\begin{array}{c|ccc}
 \text{unit}&Y_i(11,e)&Y_i(11,p)&Y_i(01)\\\hline
 1&1&3&2\\
 2&0&2&2\\
 3&2&2&1\\
 4&1&3&3
\end{array}
\]
The column means are \(1\), \(5/2\) and \(2\). In a setting where each
participant independently receives \(e\) with probability \(1/2\),
\(\tau_\lambda=(1+5/2)/2-2=-1/4\). In a setting where every participant
receives \(p\), \(\tau_\lambda=5/2-2=1/2\). The two settings share the
potential responses and the assignment law, yet the assignment comparison
estimates contrasts of opposite sign, as part~3 of
\cref{thm:p15-w02b-p15-076-versions} permits. The version-specific targets
are \(\theta_{(e)}=1-2=-1\) and \(\theta_{(p)}=1/2\). If versions are
not recorded, exchanging the two version columns for every participant
leaves every recorded response unchanged and the version law the same in
the first setting, but makes \(\theta_{(e)}=1/2\). Under the second
setting no participant receives \(e\), so \(\theta_{(e)}\) is not
determined even when versions are recorded. Participant \(3\) alone
satisfies the stable-version premise.
\end{hbexample}

\section{Interference and the assignment contrast}
\label{sec:p15-w02b-p15-076-interference}

\begin{hbdefinition}{def:p15-w02b-p15-076-interference}{The interference observation model and cluster structure}
In the standing setting, let each unit have a real response function
\(Y_i:\mathcal T^N\to\mathbb R\) and let the observation model be
\(Y_i^{\mathrm{rec}}=Y_i(Z)\). \emph{No interference} holds when
\(Y_i(z)=Y_i(z')\) whenever \(z_i=z'_i\); then \(Y_i(t)\) denotes the
common value at \(z_i=t\). For \(t\in\mathcal T\), \(t\mathbf 1\) is the
vector assigning every unit to \(t\), and the \emph{common-assignment
contrast} is \(\tau_{\rm com}=N^{-1}\sum_i\sum_tw_tY_i(t\mathbf 1)\).
The \emph{design-averaged contrast} is
\[
 \tau_\varpi=\frac1N\sum_i\sum_tw_t
   \sum_{z:\,z_i=t}\frac{\varpi(z)}{\pi_i(t)}\,Y_i(z).
\]
For a partition of the units into clusters, with \(g(i)\) the cluster of
unit \(i\), \emph{interference confined to clusters} means that
\(Y_i(z)=Y_i(z')\) whenever \(z_j=z'_j\) for every \(j\in g(i)\). A design
\emph{assigns by cluster} when \(\varpi(z)>0\) only for vectors \(z\)
constant on every cluster.
\index{interference}\index{cluster assignment}
\end{hbdefinition}

\begin{hbtheorem}{thm:p15-w02b-p15-076-interference}{proved=annex:p15-w02b-p15-076-interference}{What the assignment contrast estimates under interference}
In the setting of \cref{def:p15-w02b-p15-076-interference}:
\begin{enumerate}
\item \(\mathbb E[\widehat\tau]=\tau_\varpi\), for all response
      functions and every design law.
\item Under no interference, \(\tau_\varpi=\tau_{\rm com}=\tau\) for every
      design law.
\item With interference, \(\tau_\varpi\) depends on the design law beyond
      its marginals: there are response functions and two design laws
      with the same marginal probabilities \(\pi_i(t)\) giving different
      values of \(\tau_\varpi\), each different from \(\tau_{\rm com}\).
\item If interference is confined to clusters and the design assigns by
      cluster, then \(\tau_\varpi=\tau_{\rm com}\); the weighting statistic
      is unbiased for the common-assignment contrast.
\end{enumerate}
\end{hbtheorem}

Part~1 replaces the question ``what is the effect of cell \(t\)?'' by the
question the data answer: the mean response of a unit placed in cell \(t\),
averaged over the other units' cells as the design distributes them. A
different design for the same units, even one with the same marginal
probabilities, answers a different question. For an arbitrary interference
pattern \cref{thm:p15-w02b-p15-077-interference} shows that the
common-assignment contrast is identified exactly when each weighted common
assignment has positive design probability. Part~4 gives a restricted
interference pattern under which assigning whole clusters suffices, because
a unit's response is then fixed by its own cluster's common cell. The
price is that the cluster, not the person, becomes the independently
assigned unit; \cref{thm:p15-w02b-p15-076-families} computes that price.

\begin{hbexample}{def:p15-w02b-p15-076-interference-example}{Two participants who share a session}
Let \(N=2\), \(\mathcal T=\{0,1\}\), \(w=(-1,1)\), and let each
participant's response count the treated participants in the shared
session: \(Y_1(z)=Y_2(z)=z_1+z_2\). Then \(\tau_{\rm com}=2-0=2\).
Under complete randomization with one treated participant, the vectors
\((1,0)\) and \((0,1)\) each have probability \(1/2\). A participant in
cell \(1\) then has response \(1\), and so does a participant in cell
\(0\), so \(\tau_\varpi=0\). Under independent assignment with
probability \(1/2\) each, the marginals are again \(1/2\); a participant in
cell \(1\) has mean response \(3/2\) and one in cell \(0\) has mean
response \(1/2\), so \(\tau_\varpi=1\). Assigning the session as a whole,
with \((0,0)\) and \((1,1)\) each of probability \(1/2\), gives
\(\tau_\varpi=2=\tau_{\rm com}\), as part~4 states.
\end{hbexample}

\section{Adherence: assignment and receipt}
\label{sec:p15-w02b-p15-076-adherence}

\begin{hbdefinition}{def:p15-w02b-p15-076-adherence}{Receipt, exclusion, monotonicity and compliance types}
In the standing setting take \(\mathcal T=\{0,1\}\), \(w=(-1,1)\), and
\(0<\pi_i(1)<1\). Each unit has potential receipts \(D_i(z)\in\{0,1\}\)
and real potential responses \(Y_i(z,d)\) for \(z,d\in\{0,1\}\). The
observation model records \(D_i=D_i(Z_i)\) and
\(Y_i^{\mathrm{rec}}=Y_i(Z_i,D_i(Z_i))\). The \emph{exclusion premise}
holds when \(Y_i(0,d)=Y_i(1,d)\), written \(Y_i(d)\), for each unit and
each \(d\). \emph{Monotonicity} holds when \(D_i(1)\geq D_i(0)\) for
every unit. Units with \((D_i(1),D_i(0))=(1,0)\) are \emph{compliers},
\((0,1)\) \emph{defiers}, \((1,1)\) \emph{always-receivers} and \((0,0)\)
\emph{never-receivers}; let \(\mathcal C\) be the set of compliers. Define
the assignment contrasts
\[
 \begin{aligned}
 \mathrm{ITT}_Y&=\frac1N\sum_i\bigl(Y_i(1,D_i(1))-Y_i(0,D_i(0))\bigr),\\
 \mathrm{ITT}_D&=\frac1N\sum_i\bigl(D_i(1)-D_i(0)\bigr),
 \end{aligned}
\]
and, under exclusion, the complier contrast
\(\mathrm{CACE}=|\mathcal C|^{-1}\sum_{i\in\mathcal C}(Y_i(1)-Y_i(0))\)
when \(\mathcal C\ne\varnothing\) and the all-unit receipt contrast
\(\mathrm{ATE}_D=N^{-1}\sum_i(Y_i(1)-Y_i(0))\).
\index{adherence}\index{noncompliance}\index{compliers}\index{exclusion premise}
\end{hbdefinition}

\begin{hbtheorem}{thm:p15-w02b-p15-076-adherence}{proved=annex:p15-w02b-p15-076-adherence}{Assignment contrasts, the complier contrast and their limits}
In the setting of \cref{def:p15-w02b-p15-076-adherence}:
\begin{enumerate}
\item With no exclusion or monotonicity premise, the weighting statistics
      of \(Y^{\mathrm{rec}}\) and of \(D\) are unbiased for
      \(\mathrm{ITT}_Y\) and \(\mathrm{ITT}_D\).
\item Under exclusion,
      \(\mathrm{ITT}_Y=N^{-1}\sum_i(D_i(1)-D_i(0))(Y_i(1)-Y_i(0))\): only
      compliers and defiers contribute, with opposite signs.
\item Under exclusion and monotonicity, \(\mathrm{ITT}_D=|\mathcal C|/N\)
      and \(\mathrm{ITT}_Y=\mathrm{ITT}_D\cdot\mathrm{CACE}\) when
      \(\mathcal C\ne\varnothing\); if \(\mathcal C=\varnothing\), both
      assignment contrasts vanish. Hence
      \(\mathrm{CACE}=\mathrm{ITT}_Y/\mathrm{ITT}_D\) whenever
      \(\mathrm{ITT}_D>0\).
\item If no unit can receive the treatment when assigned to cell \(0\),
      that is \(D_i(0)=0\) for every unit, then monotonicity holds,
      \(\mathcal C=\{i:D_i(1)=1\}\), and under exclusion the ratio
      \(\mathrm{ITT}_Y/\mathrm{ITT}_D\) is the mean receipt effect among the
      units that receive the treatment when assigned to it.
\item Monotonicity without exclusion does not suffice: there are units
      satisfying monotonicity, with \(\mathrm{ITT}_D>0\) and with every
      complier's response unaffected by both assignment and receipt, for
      which \(\mathrm{ITT}_Y/\mathrm{ITT}_D=1\).
\item Exclusion without monotonicity does not suffice: there are units with
      \(Y_i(1)-Y_i(0)>0\) for every unit, \(\mathrm{ITT}_D>0\),
      \(\mathrm{ITT}_Y=0\) and \(\mathrm{CACE}=2\).
\end{enumerate}
\end{hbtheorem}

Part~1 is the formal content of an analysis by assigned cell: it estimates
the effect of assignment, which includes whatever uptake the assignment
produces, without any premise about receipt. A receipt effect requires
further premises, and parts~5 and~6 show that each of the two used in
part~3 does work: dropping either one breaks the identity. The complier
contrast concerns a set of units that is not observed: a recorded pair
\((Z_i,D_i)=(1,1)\) is compatible with a complier and with an
always-receiver. \Cref{thm:p15-w02b-p15-077-noncompliance} gives the
corresponding structural result: the effect of receipt for all units is not
identified from the recorded law even under exclusion and randomization.

\begin{hbexample}{def:p15-w02b-p15-076-adherence-example}{Six units with three compliance types}
Stipulate the following units, with exclusion holding and larger responses
better. Assign three of the six units to cell \(1\) by complete
randomization, so \(\pi_i(1)=1/2\).
\[
\begin{array}{c|cccc|c}
 \text{unit}&D_i(0)&D_i(1)&Y_i(0)&Y_i(1)&\text{type}\\\hline
 1&0&1&1&3&\text{complier}\\
 2&0&1&2&3&\text{complier}\\
 3&0&0&2&4&\text{never-receiver}\\
 4&0&0&0&0&\text{never-receiver}\\
 5&1&1&1&3&\text{always-receiver}\\
 6&0&1&1&1&\text{complier}
\end{array}
\]
Here \(\mathrm{ITT}_Y=(2+1+0)/6=1/2\), \(\mathrm{ITT}_D=1/2\), and
\(\mathrm{CACE}=(2+1+0)/3=1\), the ratio of the two assignment contrasts.
The all-unit receipt contrast is \(\mathrm{ATE}_D=7/6\). An
\emph{as-treated} comparison contrasts the recorded responses of units
that received the treatment with those of units that did not. Its
population counterpart,
\[
 \frac{\sum_i\mathbb E[D_iY_i^{\mathrm{rec}}]}{\sum_i\mathbb E[D_i]}
 -\frac{\sum_i\mathbb E[(1-D_i)Y_i^{\mathrm{rec}}]}{\sum_i\mathbb E[1-D_i]}
 =\frac{13}{5}-\frac87=\frac{51}{35},
\]
equals none of \(1/2\), \(1\) and \(7/6\): the units that receive the
treatment differ in type from those that do not.
\end{hbexample}

\section{Missing outcomes: what each premise licenses}
\label{sec:p15-w02b-p15-076-missing}

\begin{hbdefinition}{def:p15-w02b-p15-076-missing}{Potential recording, complete-case and missing means}
In the standing setting with the response observation model of
\cref{def:p15-factorial-assignment}, let each unit have potential
recording indicators \(R_i(t)\in\{0,1\}\). The recorded variables are
\(Z\), \(R_i=R_i(Z_i)\) and \(R_iY_i(Z_i)\); an unrecorded response is
not seen. Assume a known range: \(\ell\leq Y_i(t)\leq u\) for all
\(i,t\). Write \(\mu(t)=N^{-1}\sum_iY_i(t)\), and
\[
 A_t=\frac1N\sum_iR_i(t)Y_i(t),\qquad
 \bar r_t=\frac1N\sum_iR_i(t),\qquad q_t=1-\bar r_t.
\]
When \(\bar r_t>0\), the \emph{complete-case mean} is
\(\mu_{\rm cc}(t)=A_t/\bar r_t\); when \(q_t>0\), the \emph{missing
mean} is \(\mu_{\rm mis}(t)=(N q_t)^{-1}\sum_i(1-R_i(t))Y_i(t)\). The
finite-unit covariance is
\(\operatorname{Cov}_N(R(t),Y(t))=A_t-\bar r_t\mu(t)\).
For a baseline covariate \(X_i\), recorded for every unit before
assignment and taking finitely many values, the same quantities computed
over the stratum \(\{i:X_i=x\}\), of size \(N_x\), carry the argument
\((t,x)\).
\index{missing outcomes}\index{complete-case analysis}
\end{hbdefinition}

\begin{hbtheorem}{thm:p15-w02b-p15-076-missing}{proved=annex:p15-w02b-p15-076-missing}{Missing outcomes: bounds and bias}
In the setting of \cref{def:p15-w02b-p15-076-missing}:
\begin{enumerate}
\item The weighting statistics of \(\mathbf 1_{\{Z_i=t\}}R_iY_i(Z_i)\)
      and \(\mathbf 1_{\{Z_i=t\}}R_i\), each divided by \(\pi_i(t)\) and
      averaged over units, are unbiased for \(A_t\) and \(\bar r_t\).
\item \emph{No premise on recording.} The set of values of \(\mu(t)\) over
      all values in \([\ell,u]\) of the unrecorded responses, with every
      \(R_i(t)\) and every recorded response held fixed, is exactly
      \([A_t+\ell q_t,\ A_t+uq_t]\). The corresponding set for \(\tau\) is
      the interval with endpoints
      \[
       \sum_tw_tA_t+\sum_tq_t\min(w_t\ell,w_tu),\qquad
       \sum_tw_tA_t+\sum_tq_t\max(w_t\ell,w_tu),
      \]
      of width \((u-\ell)\sum_t|w_t|q_t\).
\item \emph{Complete-case bias.} If \(\bar r_t>0\), then
      \(\mu_{\rm cc}(t)-\mu(t)=\operatorname{Cov}_N(R(t),Y(t))/\bar r_t\);
      if also \(q_t>0\), this equals
      \(q_t(\mu_{\rm cc}(t)-\mu_{\rm mis}(t))\), and its absolute value is
      at most \(q_t(u-\ell)\). Hence \(\mu_{\rm cc}(t)=\mu(t)\) if and only
      if recording and response are uncorrelated across the units in cell
      \(t\).
\item \emph{Recording unrelated to the response within covariate strata.}
      If for every stratum \(x\) one has \(\bar r_t(x)>0\) and
      \(\operatorname{Cov}_N(R(t),Y(t))(x)=0\), then
      \(\mu(t)=\sum_x(N_x/N)\,\mu_{\rm cc}(t,x)\). This premise does not
      imply \(\mu_{\rm cc}(t)=\mu(t)\).
\end{enumerate}
\end{hbtheorem}

The theorem orders three premises by strength. Without any premise
about recording, the data fix \(\mu(t)\) only within an interval whose
width is the missing fraction times the range. The interval is sharp: it
cannot be narrowed without information about the unrecorded responses.
Part~3 is an exact identity, not an approximation: the complete-case mean
is off by the missing fraction times the difference between recorded and
unrecorded means. A small missing fraction therefore limits the bias only
through the range of the outcome; no fraction is small enough in itself.
Equal recording fractions in the two arms do not remove the bias either,
since the identity concerns the relation between recording and response
within each arm. Part~4 states a premise about the unrecorded responses
that licenses a covariate-weighted complete-case mean.
\Cref{thm:p15-w02b-p15-077-missingness} gives the corresponding statement
for random recording that does not depend on the response given the
assigned cell, and proves that the observed law cannot distinguish that
premise from recording that depends on the response.

\begin{hbexample}{def:p15-w02b-p15-076-missing-example}{Missing counts of inconsistent pairs}
In the design of \cref{def:p19-proposed-factorial-design} the primary
outcome counts inconsistent pairs on four tasks, so \(\ell=0\) and
\(u=4\). Take the aid-present cells with weights \(w_{01}=-1\),
\(w_{11}=1\), and suppose \(\bar r_{11}=0.8\),
\(\mu_{\rm cc}(11)=1\), \(\bar r_{01}=0.9\) and
\(\mu_{\rm cc}(01)=1.5\); these are stipulated planning values. Then
\(A_{11}=0.8\), \(A_{01}=1.35\), \(\mu(11)\in[0.8,1.6]\),
\(\mu(01)\in[1.35,1.75]\), and the aid-present contrast lies in
\([-0.95,0.25]\), an interval of width \(4\times0.3=1.2\) that contains
zero although the complete-case contrast is \(-0.5\). If the unrecorded
participants in cell \(11\) have mean \(3\), part~3 gives
\(\mu_{\rm cc}(11)-\mu(11)=0.2\times(1-3)=-0.4\), so \(\mu(11)=1.4\).

In a second stipulated scenario, let a baseline numeracy stratum split
cell \(11\)'s units equally into
low and high. Suppose the low stratum has \(\bar r_{11}(\text{low})=0.6\)
and complete-case mean \(2\), and the high stratum has
\(\bar r_{11}(\text{high})=1\) and complete-case mean \(0.6\). The
arm-level complete-case mean is
\((0.5\times0.6\times2+0.5\times0.6)/(0.5\times0.6+0.5)=1.125\). Under
the premise of part~4, \(\mu(11)=0.5\times2+0.5\times0.6=1.3\). The
complete-case mean is biased here because recording depends on the
stratum, and the stratum is related to the response.
\end{hbexample}

\begin{hbwarning}{def:p15-w02b-p15-076-design-premises}{The four premises are not checked by the assignment comparison}
A randomized assignment secures the independence of assignment from the
potential responses. It does not secure a single version, the absence of
interference, adherence or complete recording. Each of
\cref{thm:p15-w02b-p15-076-versions,thm:p15-w02b-p15-076-interference,thm:p15-w02b-p15-076-adherence,thm:p15-w02b-p15-076-missing}
leaves the weighting statistic well defined and unbiased for some
estimand. The failure in each case is that this estimand is no longer the
intended one, and nothing in the recorded data signals the substitution.
A report should therefore name the estimand that the observation model
actually delivers: a version mixture under the delivery process used, a
design-averaged contrast, an assignment contrast, or a set of values
compatible with the recorded outcomes.
\end{hbwarning}

\section{Sample-size arguments under explicit models}
\label{sec:p15-w02b-p15-076-sample-size}

A sample-size argument fixes a contrast, a mean and variance model, a
design family indexed by a size \(n\), and a criterion; it then finds the
smallest \(n\) meeting the criterion. Every number it produces is
conditional on the model and on the planning values of its unknown
parameters.

\begin{hbdefinition}{def:p15-w02b-p15-076-planning}{A planning family and its precision factor}
A \emph{planning family} is a sequence of designs indexed by an integer
\(n\geq1\) such that, for a fixed estimable contrast \(\theta=c^{\mathsf
T}\beta\) in the normal linear model of
\cref{thm:p15-factorial-linear-uncertainty} with variance \(\sigma^2\),
the least-squares estimate \(\widehat\theta_n=a_{*,n}^{\mathsf T}Y\)
satisfies \(\|a_{*,n}\|^2=\kappa_1/n\) for a constant
\(\kappa_1>0\), the \emph{precision factor}, and the fitted model leaves
\(\nu_n\geq1\) residual degrees of freedom. Thus
\(\operatorname{Var}(\widehat\theta_n)=\sigma^2\kappa_1/n\). The criterion
fixes a confidence level \(1-\eta\) and either a half-width \(h>0\), or an
effect size \(\delta\ne0\) and a power \(1-\gamma\), with
\(0<\eta,\gamma<1\). Write \(\Phi\) for the standard normal distribution
function, \(z_p=\Phi^{-1}(p)\), \(G_\nu\) for the distribution function of
the \(\chi^2_\nu\) law, and \(t_{\nu,p}\) for the Student quantile of
\cref{thm:p15-student-law}.
\index{sample size}\index{power}\index{precision factor}
\end{hbdefinition}

\begin{hblemma}{thm:p15-w02b-p15-076-families}{proved=annex:p15-w02b-p15-076-families}{Three planning families from the existing precision formulas}
\begin{enumerate}
\item \emph{Replication.} Let \(X_1\in\mathbb R^{N_1\times p}\) have rank
      \(r\), let \(c\) be estimable for \(X_1\) with least-squares vector
      \(a_{*,1}\), and let \(X_{(n)}\) stack \(n\) copies of \(X_1\).
      Then \(c\) is estimable for \(X_{(n)}\),
      \(a_{*,n}=n^{-1}(a_{*,1},\ldots,a_{*,1})\),
      \(\|a_{*,n}\|^2=\|a_{*,1}\|^2/n\) and \(\nu_n=nN_1-r\). For the
      saturated two-level full factorial in \(k\) factors and any
      factorial coefficient \(\beta_S\), \(\kappa_1=2^{-k}\) and
      \(\nu_n=2^k(n-1)\), recovering the variance \(\sigma^2/(n2^k)\) of
      \cref{thm:p15-factorial-walsh}.
\item \emph{Complete blocks.} For \(L\) treatments in \(n=B\) complete
      blocks with one observation per cell and the additive model of
      \cref{thm:p15-factorial-blocks}, a treatment contrast with weights
      \(c_t\) summing to zero has \(\kappa_1=\sum_tc_t^2\) and
      \(\nu_B=(L-1)(B-1)\).
\item \emph{Cluster assignment.} Let each of \(2K\) clusters of \(m\)
      units be assigned to one of two cells, \(K\) clusters per cell, with
      responses \(Y_{gj}=\mu_{t(g)}+U_g+e_{gj}\), where all \(U_g\) and
      \(e_{gj}\) are independent normal with mean zero and variances
      \(\sigma_U^2\) and \(\sigma_e^2\). Put
      \(\sigma^2=\sigma_U^2+\sigma_e^2\) and
      \(\rho=\sigma_U^2/\sigma^2\). Each cluster mean has variance
      \(\sigma^2(1+(m-1)\rho)/m\), and the difference of cell means has
      variance \(2\sigma^2(1+(m-1)\rho)/(mK)\). The cluster means follow
      the normal linear model of two cell means with error variance
      \(\sigma_c^2=\sigma^2(1+(m-1)\rho)/m\); in it the difference of cell
      means has \(\|a_{*,K}\|^2=2/K\). The cluster means therefore form a
      planning family in \(n=K\) with \(\kappa_1=2\), variance parameter
      \(\sigma_c^2\) and \(\nu_K=2K-2\). The same \(Km\) units per cell
      assigned individually, with independent errors of variance
      \(\sigma^2\), give the variance \(2\sigma^2/(Km)\).
\end{enumerate}
\end{hblemma}

The factor \(1+(m-1)\rho\) of part~3 is the design effect: cluster
assignment multiplies the variance of the contrast by that factor
relative to individual assignment of the same number of units. With
interference confined to clusters, part~4 of
\cref{thm:p15-w02b-p15-076-interference} makes cluster assignment the
design that targets the common-assignment contrast, so this factor is the
precision cost of that target.

\begin{hbtheorem}{thm:p15-w02b-p15-076-known-variance}{proved=annex:p15-w02b-p15-076-known-variance}{Sample size with a known variance}
In a planning family with known \(\sigma>0\):
\begin{enumerate}
\item The interval \(\widehat\theta_n\pm z_{1-\eta/2}\sigma\sqrt{\kappa_1/n}\)
      has coverage exactly \(1-\eta\), and its half-width is at most \(h\)
      if and only if
      \[
       n\geq n_h=\frac{z_{1-\eta/2}^2\,\sigma^2\kappa_1}{h^2}.
      \]
      The smallest such integer is \(\lceil n_h\rceil\).
\item For \(\delta>0\), the test rejecting \(\theta\leq\theta_0\) when
      \((\widehat\theta_n-\theta_0)/(\sigma\sqrt{\kappa_1/n})>z_{1-\eta}\)
      has rejection probability \(\Phi(d_n-z_{1-\eta})\) at
      \(\theta=\theta_0+\delta\), where \(d_n=\delta\sqrt n/(\sigma\sqrt{\kappa_1})\);
      this is at most \(\eta\) for every \(\theta\leq\theta_0\). Its power at
      \(\theta_0+\delta\) is at least \(1-\gamma\) if and only if
      \[
       n\geq\frac{(z_{1-\eta}+z_{1-\gamma})^2\sigma^2\kappa_1}{\delta^2}.
      \]
\item For \(\delta\ne0\), the two-sided test rejecting when
      \(|\widehat\theta_n-\theta_0|/(\sigma\sqrt{\kappa_1/n})>z_{1-\eta/2}\)
      has level \(\eta\) and power
      \(\Phi(|d_n|-z_{1-\eta/2})+\Phi(-|d_n|-z_{1-\eta/2})\) at
      \(\theta_0+\delta\). It suffices for power at least \(1-\gamma\) that
      \(n\geq(z_{1-\eta/2}+z_{1-\gamma})^2\sigma^2\kappa_1/\delta^2\).
\end{enumerate}
\end{hbtheorem}

\begin{hbtheorem}{thm:p15-w02b-p15-076-student}{proved=annex:p15-w02b-p15-076-student}{Sample size when the variance is estimated}
In a planning family with unknown \(\sigma\), let \(s_n^2\) be the
residual variance estimate, \(\nu=\nu_n\), and
\(H_n=t_{\nu,1-\eta/2}\,s_n\sqrt{\kappa_1/n}\) the half-width of the
exact Student interval of \cref{thm:p15-factorial-linear-uncertainty}.
\begin{enumerate}
\item \(t_{\nu,p}>z_p\) for every \(\nu\geq1\) and \(1/2<p<1\).
      Consequently the interval \(\widehat\theta_n\pm z_{1-\eta/2}s_n\sqrt{\kappa_1/n}\)
      has coverage \(2F_\nu(z_{1-\eta/2})-1<1-\eta\).
\item \(\mathbb E[H_n]\leq t_{\nu,1-\eta/2}\sigma\sqrt{\kappa_1/n}\), and
      \[
       \Pr(H_n\leq h)=G_\nu\!\left(\frac{\nu\,n\,h^2}
           {t_{\nu,1-\eta/2}^2\sigma^2\kappa_1}\right).
      \]
\item Let \(T=(\widehat\theta_n-\theta_0)/(s_n\sqrt{\kappa_1/n})\),
      \(q>0\), and \(d_n\) as in \cref{thm:p15-w02b-p15-076-known-variance}.
      At \(\theta=\theta_0+\delta\),
      \[
       \begin{aligned}
       \Pr(T>q)&=\int_0^\infty\Phi\!\left(d_n-q\sqrt{w/\nu}\right)dG_\nu(w),\\
       \Pr(T<-q)&=\int_0^\infty\Phi\!\left(-d_n-q\sqrt{w/\nu}\right)dG_\nu(w),
       \end{aligned}
      \]
      and for every \(v>0\),
      \(\Pr(T>q)\geq\Phi(d_n-qv)\,G_\nu(\nu v^2)\). At \(\delta=0\) the
      tests rejecting when \(T>t_{\nu,1-\eta}\) or when
      \(|T|>t_{\nu,1-\eta/2}\) have level exactly \(\eta\).
\end{enumerate}
\end{hbtheorem}

Part~1 shows that the known-variance formula with \(z\) understates what an
analysis with an estimated variance needs, and that inserting \(z\) into
the estimated-variance interval loses coverage. Because \(\nu_n\) grows
with \(n\), the Student criteria are solved by stepping through \(n\).
Part~2 gives two criteria for the random half-width: a bound on its mean,
and the exact probability that it meets the target. Part~3 gives the
exact rejection probability as a one-dimensional integral, and a product
bound that needs only a normal and a \(\chi^2\) distribution function.

\begin{hbmethod}{def:p15-w02b-p15-076-procedure}{Computing a sample size under an explicit model}
\begin{enumerate}
\item State the contrast, its weights and the inferential target: an
      interval half-width \(h\), or an effect \(\delta\) with power
      \(1-\gamma\), at confidence level \(1-\eta\).
\item State the mean model and the variance model, and check that the
      contrast is estimable in each candidate design
      (\cref{thm:p15-factorial-estimability}). Replication does not cure an
      alias.
\item Identify the planning family, its precision factor \(\kappa_1\) and
      its residual degrees of freedom \(\nu_n\)
      (\cref{thm:p15-w02b-p15-076-families}).
\item State the planning value of \(\sigma\), with its basis.
\item Compute the known-variance size of
      \cref{thm:p15-w02b-p15-076-known-variance}; then step \(n\) upward
      until the Student criterion of
      \cref{thm:p15-w02b-p15-076-student} chosen in advance holds.
\item Report the size together with the model, the planning value and the
      power or coverage at other plausible values of \(\sigma\).
\end{enumerate}
\end{hbmethod}

\begin{hbexample}{def:p15-w02b-p15-076-factorial-size}{Replicates for a three-factor full factorial}
Plan replicates \(n\) of the three-factor full factorial so that each
factorial coefficient \(\beta_S\) has a \(95\%\) interval half-width at
most \(h=0.25\). By \cref{thm:p15-w02b-p15-076-families},
\(\kappa_1=1/8\) and \(\nu_n=8(n-1)\). Take the planning value
\(\sigma^2=5/8\), the pooled estimate of the replicated example in
\cref{sec:p15-factorial-full}. With a known variance,
\(n_h=1.959964^2\times(5/8)/(8\times0.0625)=4.8018\), so \(n=5\). With an
estimated variance, the mean criterion of part~2 of
\cref{thm:p15-w02b-p15-076-student} gives, at \(n=5\) and \(\nu=32\),
\(t_{32,0.975}\sqrt{(5/8)/40}=2.036933\times0.125=0.2546>0.25\); at
\(n=6\) and \(\nu=40\), \(2.021075\times\sqrt{(5/8)/48}=0.2306\). Thus
\(n=6\), or \(48\) runs. The probability that the realized half-width
meets the target is \(G_{40}(47.0)=0.7925\) at \(n=6\) and \(0.9603\) at
\(n=7\). The planning value is itself an estimate on \(8\) degrees of
freedom: since \(8s^2/\sigma^2\) has the \(\chi^2_8\) law, it falls below
the true variance with probability \(G_8(8)=0.5665\). If the realized
analysis used \(z\) with only those \(8\) degrees of freedom, its nominal
\(95\%\) interval would have coverage \(2F_8(1.959964)-1=0.9143\).
\end{hbexample}

\begin{hbexample}{def:p15-w02b-p15-076-block-size}{Blocks for a power target}
For \(L=4\) dosages in complete blocks, plan to detect the difference
\(\delta=1\) between dosages \(1\) and \(2\), with weights \((-1,1,0,0)\),
by a two-sided test at \(\eta=0.05\) with power \(0.8\), at planning value
\(\sigma=1\). Here \(\kappa_1=2\) and \(\nu_B=3(B-1)\). The known-variance
bound is \((1.959964+0.841621)^2\times2=15.698\), so \(B=16\), with
two-sided normal power \(0.8074\). For the Student test, numerical
evaluation of the exact integrals of part~3 of
\cref{thm:p15-w02b-p15-076-student} gives power \(0.7904\) at \(B=16\)
(\(\nu=45\), \(t_{45,0.975}=2.014103\)) and \(0.8151\) at \(B=17\)
(\(\nu=48\)); the product bound reaches \(0.8\) only at \(B=22\), where
\(v=1.177\) gives \(0.8133\). If the true standard deviation is \(1.25\)
rather than \(1\), the exact Student power at \(B=17\) falls to
\(0.6276\).
\end{hbexample}

\begin{hbexample}{def:p15-w02b-p15-076-cluster-size}{Sessions as clusters}
Suppose the participants of \cref{def:p19-proposed-factorial-design} are
tested in sessions of \(m=6\) and whole sessions are assigned, so that
interference within a session does not alter the target. For the
aid-present contrast between two cells, take planning values \(\sigma=1\),
\(\rho=0.05\) and \(\delta=0.5\), with a two-sided test at \(\eta=0.05\)
and power \(0.8\). The design effect is \(1+5\times0.05=1.25\), so
\(\sigma_c^2=1.25/6\), and part~3 of
\cref{thm:p15-w02b-p15-076-known-variance} with \(\kappa_1=2\) requires
\(K\geq(2.801585)^2\times2\times(1.25/6)/0.25=13.08\), so \(14\)
sessions, or \(84\) participants, per cell. Individual assignment of
\(n\) participants per cell with independent responses has variance
\(2\sigma^2/n\) and requires \(n\geq(2.801585)^2\times2/0.25=62.79\),
so \(63\) participants per cell. The planning values are premises, not estimates from a conducted
study.
\end{hbexample}

\begin{hbwarning}{def:p15-w02b-p15-076-size-limits}{What a sample-size argument does not establish}
\begin{enumerate}
\item The size is computed at a planning value of \(\sigma\). The
      half-width scales with the true \(\sigma\), and the power at a fixed
      \(\delta\) decreases as \(\sigma\) increases; the guarantees hold at
      the planning value, not at the true one. A planning value estimated
      from earlier data is random: the guarantee becomes a statement about
      a size that is itself a random variable, with no fixed coverage or
      power.
\item With an estimated variance the half-width is random. A size
      satisfying the mean criterion of part~2 of
      \cref{thm:p15-w02b-p15-076-student} may still give a half-width
      above \(h\) with substantial probability; the probability criterion
      has to be chosen explicitly.
\item The quantile must match the analysis: a size computed with \(z\) and
      an analysis with \(t\) give less power than planned, and an interval
      computed with \(z\) and an estimated variance gives less coverage
      than stated.
\item Power at \(\delta\) says nothing about effects smaller than
      \(\delta\), about effects of other contrasts, or about simultaneous
      error rates across many contrasts.
\item The model premises remain premises: an alias is not removed by
      replication, the clustered unit of assignment governs the variance
      however many measurements each unit supplies, and the estimand is
      the one delivered by the observation model of the earlier sections.
      A precise estimate of a version mixture or of a complete-case mean
      remains an estimate of that quantity.
\item A size computed before the data says nothing about whether a
      particular realized interval contains the target.
\end{enumerate}
\end{hbwarning}

\section{Limitations and relationships}
\label{sec:p15-w02b-p15-076-relationships}

The results are finite-unit, design-based statements about recruited units,
except for the sample-size results, which assume a normal linear model.
Generalization beyond the recruited units requires a sampling or transport
argument, as in \cref{sec:p15-w02b-p15-075-sampling}. The version and
interference results treat the delivery process and the response functions
as fixed features of the setting; when either changes, the estimand changes
with it. The adherence results concern two cells and binary receipt; a
receipt that varies in dose or timing is a sequential policy, for which
\cref{thm:p16-sequential-joint-g-formula,thm:p16-sequential-failures} give
the identification conditions and their failures, including positivity.
The missing-outcome bounds use only the outcome range; they are sharp for
that information but are often too wide to decide a sign.

The structural forms of the interference, noncompliance and missingness
boundaries are in
\cref{thm:p15-w02b-p15-077-interference,thm:p15-w02b-p15-077-noncompliance,thm:p15-w02b-p15-077-missingness};
the present chapter adds what the weighting statistic estimates when each
premise fails, the version premise, the complier identity, sharp bounds and
the exact complete-case bias. The variance formulas used for planning are
those of
\cref{thm:p15-factorial-linear-uncertainty,thm:p15-factorial-walsh,thm:p15-factorial-blocks,thm:p15-factorial-replication},
and the Student law is that of \cref{thm:p15-student-law}. The randomization
variance of the weighting statistic, which does not need a normal model,
is in \cref{thm:p15-factorial-assignment}.
